%=====================================================================
% REFERENCIAS DO MATERIAL DE LIGACAO, MONTAGEM E VERIFICACAO
%=====================================================================

\section{Referências e continuidade do estudo}

\begin{frame}{Referências normativas e fontes oficiais}
\small
\begin{itemize}
\item \textbf{ABNT NBR 5410} — instalações elétricas de baixa tensão; confirmar edição vigente no \href{https://www.abntcatalogo.com.br/}{Catálogo ABNT}.
\item \textbf{NR-10} — segurança em instalações e serviços em eletricidade: \href{https://www.gov.br/trabalho-e-emprego/pt-br/acesso-a-informacao/participacao-social/conselhos-e-orgaos-colegiados/comissao-tripartite-partitaria-permanente/normas-regulamentadora/normas-regulamentadoras-vigentes/norma-regulamentadora-no-10-nr-10}{página oficial do MTE}.
\item \textbf{Equatorial Energia Maranhão} — \href{https://ma.equatorialenergia.com.br/institucional/leis-e-normas-tecnicas/normas-tecnicas/instalacao-de-padrao/}{normas técnicas e desenhos vigentes de fornecimento}.
\item \textbf{CBMMA} — \href{https://cbm.ssp.ma.gov.br/unidades-bm/unidades-administrativas/dat/legislacao-tecnica-dat/}{legislação técnica de segurança contra incêndio}, incluindo NT 18 para iluminação de emergência.
\item Fichas técnicas, diagramas, torques, compatibilidade e procedimentos dos \textbf{fabricantes efetivamente especificados}.
\item Procedimentos, prontuário, análise de risco e regras de liberação da \textbf{organização responsável pela instalação}.
\end{itemize}
\begin{caixaAt}
Normas e manuais mudam. Registrar número, edição/revisão e data de consulta no memorial e no plano de ensaios.
\end{caixaAt}
\end{frame}

\begin{frame}{Nota temporal sobre a NR-10}
\small
Em agosto de 2026, a página oficial do Ministério do Trabalho e Emprego informa:
\begin{itemize}
\item texto vigente da NR-10 até \textbf{31 de maio de 2027};
\item nova redação aprovada pela \textbf{Portaria MTE nº 737/2026}, com vigência indicada a partir de \textbf{1º de junho de 2027}.
\end{itemize}
\begin{caixaNorma}[title=Aplicação no curso]
Usar o texto legal vigente na data da atividade e preparar a atualização do material/procedimentos antes da mudança de vigência. A existência de uma futura redação não autoriza antecipar ou ignorar obrigações atuais.
\end{caixaNorma}
\begin{caixaAt}
Este Beamer ensina leitura e método, mas não habilita profissional nem substitui capacitação, autorização, análise de risco, supervisão e procedimentos exigidos para o trabalho real.
\end{caixaAt}
\end{frame}


